\title{Large Language Models Can Automatically Engineer Features for Few-Shot Tabular Learning}

\begin{abstract}
Large Language Models (LLMs) have demonstrated an exceptional capacity to solve novel and complex reasoning tasks, suggesting significant opportunities for their application in tabular learning—a domain integral to numerous practical uses. In this study, we introduce an innovative in-context learning approach, \textsf{FeatLLM}, which leverages LLMs to perform feature engineering, thereby generating transformed datasets that are better optimized for tabular prediction problems. These newly constructed features facilitate class probability estimation through straightforward downstream machine learning models, such as linear regression, resulting in strong few-shot learning performance. At inference, \textsf{FeatLLM} relies solely on the simple predictive model equipped with the crafted features, avoiding the need to re-engage the LLM during prediction. Relative to previous LLM-based frameworks, \textsf{FeatLLM} obviates the requirement to query the LLM on each test instance. It also operates efficiently with just API-level access, effectively addressing prompt size constraints. Empirical results from diverse tabular datasets indicate that \textsf{FeatLLM} discovers superior rules, achieving average improvements of 10\% over alternative methods like TabLLM and STUNT.
\end{abstract}

\section{Introduction}
Over the past few years, Large Language Models (LLMs) have exhibited remarkable proficiency in addressing unfamiliar tasks via response generation, and do so without the necessity for domain-specific fine-tuning~\cite{creswell2022selection,imani2023mathprompter,taylor2022galactica}. The extensive pre-training on broad textual datasets allows LLMs to substitute for traditional expert knowledge across multiple domains~\cite{Genomes2021,singhal2023large,wilcox2003role}, facilitating automation in tasks such as conversational analysis~\cite{finch2023leveraging} and code correction~\cite{fan2023automated}. By supplying a combination of descriptive prompts and a limited set of examples, LLMs are able to infer task context and focus on salient features and training cases~\cite{brown2020language,zhao2023survey}, pointing to their ability to act as adept feature engineers.

Efforts to harness LLMs’ background knowledge and logical reasoning have been increasingly extended to tabular learning scenarios. For example, leveraging insights about elevated disease risk with increasing age can enhance predictive accuracy in medical contexts. In these approaches, prior studies represent tasks and feature attributes as natural language text, serialize the input data accordingly, and submit these prompts to LLMs~\cite{dinh2022lift,wang2023anypredict}. Subsequent procedures often entail in-context learning~\cite{nam2023semi} or the use of parameter-efficient fine-tuning strategies~\cite{dinh2022lift,hegselmann2023tabllm}. The integration of LLM-derived prior knowledge has been validated to substantially outperform traditional tabular models, particularly under low-data conditions~\cite{hegselmann2023tabllm}.

Existing approaches to LLM-based tabular learning encounter notable drawbacks. For instance, performing end-to-end inference necessitates at least one LLM call per instance, leading to significant computational costs. Additionally, achieving competitive predictive accuracy with most of these methods often requires fine-tuning the language models. This prerequisite renders their deployment impractical on LLMs lacking available training code or tuning APIs—a particularly pressing issue as many state-of-the-art LLMs provide only restricted access through inference APIs~\cite{anil2023palm,openai2023gpt4}. Moreover, the majority of these strategies are ill-equipped to handle extensive prompts~\cite{liu2023lost}; as the number of features in tabular datasets increases and the input sequence length grows, these methods either become unmanageable or experience a drop in predictive performance. Collectively, these issues limit the feasibility and scalability of applying LLM-based tabular modeling in real-world settings.

The methodology we propose diverges from the prevailing paradigm of employing LLMs solely for end-to-end prediction tasks. Rather, we harness LLMs primarily for feature construction, capitalizing on their embedded world knowledge for more judicious feature extraction. We present an innovative in-context learning strategy, \textsf{FeatLLM}, that seeks to elucidate the "criteria" that drive LLM decisions. For example, in medical diagnosis tasks, the LLM can articulate explicit rules specifying which feature values are indicative of a particular disease. These extracted criteria serve as the foundation for new, informative features that can be substituted for the original set in downstream analyses. This shift repositions the LLM as a feature generator, thus simplifying the subsequent modeling stage—once salient features have been identified, downstream inference no longer demands complex machine learning models, resulting in reduced inference time. Through the in-context learning capability of LLMs, this process can occur with zero parameter updates, leveraging prompt engineering and examples to uncover relevant features. Additionally, integrating ensemble learning strategies such as feature bagging~\cite{breiman2001random} allows for circumventing issues posed by lengthy prompts, further enhancing the approach's scalability.

\textsf{FeatLLM} decomposes the process of engineering novel features from input prompts into two primary sub-tasks: (1) comprehending the problem and deducing the connections between input features and target variables; and (2) using insights derived from the initial step, along with few-shot examples, to formulate discriminative rules for class separation. This sequential reasoning strategy guides LLMs to disregard non-informative features in the process of rule generation. The resulting rules are executed on the dataset (by parsing them as code), thereby converting the original samples into binary indicators that reflect compliance with each respective rule. Subsequently, a simple model is trained to estimate class membership probabilities using these binary features. To further enhance stability, an ensemble technique is adopted, wherein feature generation is conducted multiple times alongside feature bagging. By strategically tuning the number of features and data points utilized in each bag, \textsf{FeatLLM} effectively addresses the issue of oversized prompts. Notably, this framework avoids the need for additional training and operates with minimal inference expense, making it suitable for diverse real-world tasks.

The effectiveness of \textsf{FeatLLM} is demonstrated on 13 tabular datasets under limited data availability, showcasing its consistently strong performance. The results highlight the ability of LLMs to leverage both existing and data-driven knowledge to derive meaningful features. Compared to recent few-shot learning baselines, our approach attains superior results under a range of experimental scenarios. Implementation details and resources are provided through the anonymized GitHub repository at~\texttt{https://github.com/Sungwon-Han/FeatLLM}.

\section{Related Work}

\subsection{Few-Shot Learning with Tabular Data} Progress in few-shot learning has allowed models to develop transferable representations using only a small number of labeled samples, substantially lowering annotation burdens~\cite{chen2018closer}. Although the field initially concentrated on images, its robust capacity to generalize has also been demonstrated in other modalities—such as text, audio, and increasingly, tabular data~\cite{majumder2022few,nam2023stunt,schick2021exploiting}. Nonetheless, working with scant labeled examples can heighten the chance of models capturing misleading or coincidental associations~\cite{bartlett2021deep}. To mitigate this, recent research has investigated the use of supplementary information sources and the integration of prior, domain-relevant knowledge, thereby introducing beneficial inductive biases during training. This includes leveraging large pools of unlabeled data alongside strategic data augmentation to support semi-supervised methods~\cite{ucar2021subtab,yoon2020vime}, or acquiring strong model initialization via broad, task-agnostic unsupervised pre-training~\cite{somepalli2021saint}. In unsupervised settings, models may use objectives based on masking or corrupting input data followed by reconstruction~\cite{arik2021tabnet,majmundar2022met}, or rely on contrastive learning tasks grounded in data augmentations~\cite{bahri2022scarf,chen2020simple}. However, the inherent heterogeneity of tabular datasets complicates the creation of generic augmentation techniques, which has limited the success of such methods in few-shot scenarios~\cite{nam2023stunt}.

Emerging strategies involve pre-training models on a wide spectrum of tabular datasets, rather than restricting training to data closely resembling the target application, and then transferring this knowledge to downstream tasks~\cite{zhang2023generative,levin2022transfer,zhu2023xtab}. An illustrative example is TransTab~\cite{wang2022transtab}, which applies transformer models to learn column-level semantic relationships through column names in addition to cell values. This approach enables the extraction of generalized knowledge from disparate tables and columns, supporting robust zero-shot and few-shot task inference. Similarly, TabPFN~\cite{hollmann2022tabpfn} employs a variety of distributions to synthesize data used for pre-training a Bayesian neural network; subsequent inference incorporates both training and test sets from the target domain, removing the need for continued adaptation. Exposure to diversified data distributions enables these models to consistently outperform standard tree-based algorithms, offering distinct advantages in low-data contexts.

\subsection{Language-Interfaced Tabular Learning} Utilizing prior knowledge via large language models (LLMs)~\cite{anil2023palm,openai2023gpt4} represents another fruitful pathway. Thanks to their training on massive and heterogeneous text datasets, LLMs can adapt to previously unseen tasks and demonstrate strong cross-domain applicability~\cite{brown2020language}. Recent explorations focus on tapping into the extensive knowledge embedded within LLMs for tackling tabular problems. This typically entails formatting tables as text sequences and embedding them into LLM input prompts~\cite{dinh2022lift,hegselmann2023tabllm,wang2023anypredict}. Input prompts are augmented with tabular entries, feature metadata, and task annotations, allowing the LLM to infer relationships between variables and predict outcomes. The state-of-the-art is moving towards two main directions: either adapting model parameters with lightweight techniques such as LoRA~\cite{hu2021lora} and IA3~\cite{liu2022few}, or leveraging in-context learning, where the LLM is provided with few-shot exemplars directly within the prompt, eliminating the need for further training~\cite{dinh2022lift,hegselmann2023tabllm,nam2023semi,slack2023tablet}.

Although the approaches discussed above significantly enhance few-shot learning on tabular data, the necessity of routing all inference through the LLM restricts their practical use in industrial settings. This paper proposes to move beyond the standard black-box paradigm, wherein the LLM performs end-to-end prediction for every data point. Rather, our method leverages the LLM to explicitly deduce the characteristic rules or decision criteria for each output class. \textsf{FeatLLM} translates these generated rules into executable code for feature construction, thus facilitating inference that circumvents the LLM. This strategy takes advantage of in-context learning to draw on both existing knowledge and limited sample data when deriving the governing rules.

\section{Methods}
\textsf{FeatLLM} is built to identify and extract the defining ``rules" that specify membership to each output class, using a small set of example samples as illustrated in Figure~\ref{fig:main_model}. Through interactions with the LLM, these inferred rules are converted into binary indicators for each input sample, signifying the presence or absence of rule satisfaction. The resulting binary feature matrix is then linearly combined—using dot products with non-negative, trainable coefficients—to compute a class likelihood score. The process of training the model enables the estimation of each rule’s relative impact on class assignment (see Section~\ref{Sec:training} for details). To improve robustness, this sequence is repeated over multiple bootstrap samples, with final predictions aggregated through an ensemble approach. The following subsections present a detailed discussion of the problem definition and a breakdown of each phase of the methodology.

\vspace*{-0.12in}
\paragraph{Problem formulation.} Consider a tabular dataset $\mathcal{D} = \{(\mathbf{x}^i, \mathbf{y}^i)\}_{i=1}^{N}$, comprising $N$ annotated instances. Each data point $\mathbf{x}^i$ is a vector in $\mathbb{R}^d$, representing $d$ individual input features, while each corresponding label $\mathbf{y}^i$ belongs to a fixed set of classes, as typical in classification tasks. The dataset includes natural language feature names, collected in $F = \{f_j\}_{j=1}^{d}$, and their definitions, which are supplied separately. For few-shot (specifically $k$-shot, with $k < N$) scenarios, a subset of $k$ randomly selected, labeled samples is used for model training.

\subsection{Prompt Design for Extracting Rules}
\label{Sec:prompt_design}
To improve the accuracy of rule extraction by the LLM and emulate the logical process an expert might follow in analyzing tabular data, we structure the prompt to encourage stepwise problem solving. As outlined in Fig.~\ref{fig:prompt_for_rule} and discussed further in Appendix~\ref{sec:example_rule}, the prompt is organized into three primary segments:

\vspace*{-0.12in}
\paragraph{Basic information description.} This section conveys the foundational details necessary to address the task (highlighted in orange in Fig.~\ref{fig:prompt_for_rule}), including a statement outlining the problem with explicit label criteria, contextual explanations for each feature, and illustrative examples. The task is posed as an explicit question (such as, ``Does this patient's data regarding myocardial infarction complications suggest chronic heart failure? Yes or no?"). Further instances are detailed in Appendix~\ref{sec:task_desc}. Each feature is briefly described by its type and may be supplemented with a definition; categorical variables may be accompanied by sample category values. In the example demonstrations, several selected training examples are converted into textual format along with their annotated labels, following serialization protocols adopted in earlier research~\cite{dinh2022lift,hegselmann2023tabllm}:

\begin{align}
&\text{Serialize}(\mathbf{x}^i,\mathbf{y}^i, F) = \nonumber \\ 
&``\ f_1\ \text{is}\ \mathbf{x}^i_1.\ ... \ f_d\  \text{is}\  \mathbf{x}^i_d.\ \ \text{Answer:}\  \mathbf{y}^i\ ”
\end{align}

\vspace*{-0.12in}
\paragraph{Reasoning instruction.} Our objective is to bolster the reasoning capacity of the LLM by explicitly supplying reasoning instructions (illustrated in blue in Fig.~\ref{fig:prompt_for_rule}). These instructions initiate with a statement reminiscent of the chain-of-thought methodology~\cite{wei2022chain}, and subsequently partition the model's inferential process into two distinct phases. Initially, the LLM is tasked with independently reasoning about the causal effects or tendencies connecting the input features to the problem description, abstaining from referencing any sample demonstrations. This encourages the LLM to utilize its own general knowledge or common sense, thereby drawing upon its underlying priors relevant to the task. In the second phase, the model incorporates insights from the preceding step together with provided demonstration examples, thereby formulating explicit rules for each category. Such a bifurcated strategy mitigates the risk of the model drawing upon superficial correlations from unrelated features and instead directs its focus to more salient characteristics. To counterbalance the dangers of extracting either too few rules (underfitting) or an unwieldy number (resulting in bloated prompts), we designate a constant rule count (set at 10) for each class\footnote{Refer to Section~\ref{sec:analysis} for an examination of rule quantity choices.}. Additionally, when constructing the rules, we instruct the LLM to consider the data type of each feature—detailed in the initial information—to prevent the generation of rules with incompatible feature types.

\vspace*{-0.12in}
\paragraph{Response instruction.} To promote consistency and facilitate downstream processing of the extracted rules, we guide the LLM's outputs using tailored response instructions (see yellow markings in Fig.~\ref{fig:prompt_for_rule}). These guidelines specify both the general organizational format expected in responses (format for the response) and the structure each individual rule should adhere to (format for the rule) across different classes. By providing these detailed templates, we enable the LLM to align its output format to the type of feature it is reasoning about. Absent any special constraints, the LLM is permitted to connect multiple rules using logical connectors such as AND and OR. An illustration of the finalized outputs is available in Appendix~\ref{sec:outcome-rule}.

\subsection{Inferring Class Likelihood via Rules}
\label{Sec:training}

\paragraph{Parsing rules for feature generation.} In this stage, the previously identified rules are utilized to derive new binary features. Each class is assigned a distinct feature that reflects whether a given instance complies with that class's rule set—this is encoded as '0' or '1'. These binary indicators provide useful input for predicting class membership. As the rules produced by the LLM are expressed in natural language, an effective parsing method is essential to automate feature construction. Although efforts are made to standardize response formats and constrain the syntax during rule extraction, outputs from the LLM may still include various minor syntactic inconsistencies~\cite{kovavc2023large}, such as using verbal alternatives to mathematical operators (for example, writing ``greater than'' instead of ``$>$''). This inherent variability poses extra challenges to developing reliable parsers.

To overcome these obstacles related to processing noisy textual data, rather than engineering intricate parsing algorithms manually, we utilize the LLM itself for parsing tasks. Due to its proficiency in interpreting intended meanings even in the presence of syntax variation, and its training on both natural language and code, the LLM is well-suited for generating straightforward code from instructions and descriptive prompts. We exploit this by structuring prompts with relevant details, including the desired function name, specifications of inputs and outputs, and the extracted rules, before presenting them to the LLM. Illustrative examples of such prompts and their corresponding generated results are presented in Figures~\ref{fig:prompt_for_parsing} and ~\ref{sec:example_parsing}-\ref{sec:example_parse_outcome} in the Appendix. Finally, the code synthesized by the LLM is executed with Python’s \texttt{exec()} method to carry out the feature transformation process.

\vspace*{-0.12in}
\paragraph{Inferring class likelihood.} When new binary features are produced by rules corresponding to each class, a basic approach for assessing the likelihood that a sample belongs to each class is to tally the number of satisfied rules in each class—equivalent to summing the class-specific binary features. Nevertheless, the influence of each rule and associated feature can differ, so their relative importance must be learned from the training data. To capture these importance weights, we employ a standard machine learning (ML) technique, which operates separately over the binary features of each class. For illustration, consider a linear model that omits a bias term. The binary features for a sample $\mathbf{x}^i$ and class $k$ are expressed as $\mathbf{z}_k^i \in \{0, 1\}^R$, with $R$ indicating the number of rules assigned per class, and $c$ the total number of classes. The prediction probability vector $\mathbf{p}^i$ for each class is determined by:

\begin{align}
\text{logit}_k^i &= \max(\mathbf{w}_k, 0) \cdot \mathbf{z}_k^i, \nonumber \\
\mathbf{p}^i &= \text{Softmax}({[\text{logit}_{1}^i, ..., \text{logit}_{c}^i]}). \label{eq:infer_prob}
\end{align}

In this setting, $\mathbf{w}_k$ denotes the set of trainable weights in the linear model, encoding the relative contributions of individual features. Restricting the weights to be nonnegative by projecting onto the positive subspace ensures that, during training, features produced by the LLM cannot diminish a class's likelihood. Subsequently, the computed logits are transformed into probability estimates over classes using the softmax function. The weights $\mathbf{w}_k$ are optimized by minimizing the cross-entropy loss on a limited set of training samples.

\vspace*{-0.12in}
\paragraph{Ensembling with bagging.} In the final stage, the entire method is performed iteratively to generate several separate inference models, which are then combined via ensemble for robust final prediction. Specifically, after obtaining trained weights from $T$ independent runs, denoted as $\{\mathbf{w}[t]\}_{t=1}^T$, the class probabilities for each sample $\mathbf{p}^i$ are averaged across all sets of weights, as computed by Eq.~\ref{eq:infer_prob}, to arrive at the final decision.

To inject stochasticity into the rule generation process during ensemble construction, we employ a relatively elevated temperature setting for LLM inference. Additionally, the sequence of few-shot demonstrations is randomized, motivated by prior findings that permutation in demonstration order can influence LLM output~\cite{lu2022fantastically}\footnote{Analysis on the rule diversity can be found in Appendix~\ref{sec:diversity}}. To exploit predictive variability for more resilient predictions, we utilize bagging, whereby each iteration samples a subset of features or examples—this approach is particularly beneficial when the dataset comprises a large number of instances or attributes. The ensemble methodology proposed herein provides two key benefits. First, the inherent variability means that even when certain generated rule sets are incorrect, other, more accurate rule sets from separate trials can provide compensatory effects, thereby strengthening the robustness of the framework. This is aligned with the principle of LLM self-consistency~\cite{wang2022self}, whereby rules that appear in several independent trials tend to be more reliable. Second, the ensemble mechanism mitigates the constraint posed by the LLM's finite prompt size; bagging can be used to partition the data, reducing prompt length and maintaining stable model behavior. Whereas a single trial’s rules may incompletely represent all relevant feature relationships, an ensemble constructed from multiple weak learners across numerous trials can offset this partial coverage, minimizing any adverse effect from limited feature or instance exposure.

\section{Experiments}
We assess the effectiveness of \textsf{FeatLLM} on a suite of diverse tabular datasets and undertake a variety of analytical studies to elucidate the operational characteristics of our framework. Further experimental results—including studies involving alternative LLMs, qualitative insights, and more—are available in the Appendix due to page limitations.

\subsection{Performance Evaluation}
\paragraph{Datasets.}
The empirical studies leverage 13 benchmark datasets pertinent to binary and multi-class classification tasks: (1) Adult~\cite{asuncion2007uci}, used to predict whether an individual's income surpasses \$50,000 per year; (2) Bank~\cite{moro2014data}, which predicts a client’s likelihood of subscribing to a term deposit; (3) Blood~\cite{yeh2009knowledge}, for anticipating repeat blood donations from individuals; (4) Car~\cite{kadra2021well}, for classifying automobile quality; (5) Communities~\cite{misc_communities_and_crime_183}, for forecasting crime rates in a geographical region; (6) Credit-g~\cite{kadra2021well}, for discerning whether a person presents a good or bad credit risk; (7) Diabetes\footnote{\texttt{kaggle.com/datasets/uciml/pima-indians-diabetes-database}}, tasked with diagnosing diabetes; (8) Heart\footnote{\texttt{kaggle.com/datasets/fedesoriano/heart-failure-prediction}}, which predicts the risk of coronary artery disease; and (9) Myocardial~\cite{misc_myocardial_infarction_complications_579}, designed to determine the presence of chronic heart failure.

We further expand our evaluation by integrating newer and synthetic datasets that have not been widely examined nor included in LLM pre-training data. These include: (10) Cultivars, tasked with forecasting the potential grain yield of specific soybean cultivars\footnote{\texttt{archive.ics.uci.edu/dataset/913}}; (11) NHANES, which involves predicting the age group of an individual from their health record\footnote{\texttt{archive.ics.uci.edu/dataset/887}}; (12) Sequence-type, wherein a sequence is categorized as arithmetic, geometric, Fibonacci, or Collatz; and (13) Solution-mix, a binary task to determine if the concentration percentage of a mixture formed by four solutions surpasses 0.5. The first two datasets were added to the UCI Machine Learning Repository after September 2023, guaranteeing their absence from the training data of the LLM API (gpt-3.5-turbo-0613) underpinning our model. The last two are constructed synthetically, thereby eliminating the risk of being memorized by the LLM API, and were specifically introduced for evaluation purposes.

The spectrum of dataset sizes and complexities is summarized in Table~\ref{Tab:basic-desc}. Each dataset is equipped with explicit attribute names and descriptions. Datasets such as `Communities' and `Myocardial,' each exceeding 100 attributes, introduce considerable complexity given the restricted context window size of current LLMs.

\vspace*{-0.12in}
\paragraph{Baselines.}
Our experimental comparisons involve nine distinct baselines. The initial set of three represents standard methods for learning from tabular data: (1) Logistic regression (LogReg), (2) XGBoost~\cite{chen2016xgboost}, and (3) RandomForest~\cite{ho1995random}. The following pair are based on the assumption of having access to unlabeled data: (4) SCARF~\cite{bahri2022scarf} and (5) STUNT~\cite{nam2023stunt}, though in practice, acquiring ample quantities of unlabeled data can be impractical. Additionally, (6) TabPFN~\cite{hollmann2022tabpfn} provides a recent alternative by pretraining on a broad array of synthetic datasets. The remaining three baselines utilize LLMs: (7) In-context learning (In-context)~\cite{wei2022emergent}, (8) TABLET~\cite{slack2023tablet}, and (9) TabLLM~\cite{hegselmann2023tabllm}. In-context learning injects a handful of training examples directly into the prompt, foregoing further tuning. TABLET leverages guidance from an external classifier by adding rule-based hints and prototype information to the prompt to bolster inference. TabLLM makes use of parameter-efficient fine-tuning such as IA3~\cite{liu2022few} to adapt the LLM to the few-shot data.

\vspace*{-0.12in}
\paragraph{Implementation details.}
\textsf{FeatLLM} adopts GPT-3.5 as its default LLM backbone; however, its architecture is intended to accommodate various LLMs (see Appendix~\ref{sec:backbones} for evaluations using alternative backbones). The LLM's inference temperature is fixed at 0.5, while top-p remains set at the API’s standard value of 1. We configure the ensemble size and the rule extraction count to 20 and 10, respectively. Further analysis on hyper-parameter selection is presented in Figure~\ref{fig:hyperparameter2} and detailed in Appendix~\ref{sec:hyper-parameter}. For optimizing the linear model, Adam is employed with a learning rate of 0.01 across 200 training epochs. Optimal epochs are identified through $k$-fold cross-validation. In baseline reproductions, GPT-3.5 is used for in-context learning tasks, whereas TabLLM relies on T0~\cite{sanh2022multitask} according to its original configuration. Notably, TabLLM restricts usage to LLMs with available public checkpoints, precluding access to GPT-3.5 in such experiments. For traditional models such as LogReg, XGBoost, and RandomForest, hyper-parameters are tuned through a Grid-search procedure coupled with $k$-fold cross-validation. Here, $k$ is chosen as either 2 or 4 to ensure each class is present at least once in training splits. Expanded implementation specifics are available in Appendix~\ref{sec:baselines}.

\vspace*{-0.12in}
\paragraph{Results.}
Tables~\ref{tab:main1} and \ref{tab:main2} provide comparative performance evaluations across multiple datasets. Each experiment was carried out three times using different random splits of training and test sets, with the results summarized as means and standard deviations of AUC scores. \textsf{FeatLLM} persistently achieves either the highest or second-highest performance relative to other baseline approaches. Notably, our method can reach levels of accuracy on par with standard tabular baselines that are supplied with 32 shots, even when only 4 shots are available (refer to Fig.~\ref{fig:compares}). Although both STUNT and TabPFN achieve notable gains on certain datasets, their overall advantage compared to standard machine learning approaches remains limited. Furthermore, LLM-based models—including in-context learning, TABLET, and TabLLM—were unable to conduct inference on tabular datasets characterized by a large feature set. Our proposed method, integrating both bagging and ensemble strategies, remains effective on high-dimensional datasets and delivers robust performance. Importantly, our bagging and ensemble strategies could be incorporated into other LLM-driven frameworks; however, such incorporation would effectively double the per-sample inference cost, thus greatly elevating computational overhead and diminishing practicality.

\subsection{Analysis \& Discussion}
\label{sec:analysis}
\paragraph{Ablation study.} To assess the contributions of key components to model effectiveness, we conducted ablation studies that selectively remove or modify core aspects: (1) \textsf{-Tuning}: excludes the linear model’s weight adjustment step, using a sum of binary features for logit computation instead; (2) \textsf{-Ensemble}: disregards the ensemble process; (3) \textsf{-Description}: omits feature description input; and (4) \textsf{-Reasoning}: eliminates Step-1 in reasoning instructions during rule generation.

Table~\ref{tab:ablation} presents average AUC changes across all datasets resulting from each ablation. In every scenario, altering any component leads to a decrease in performance. Notably, the advantage conferred by weight tuning increases as the number of shots rises, suggesting that, with larger datasets, more accurate identification of rule significance becomes possible and, consequently, the impact of tuning is amplified. Conversely, the influences of feature descriptions and reasoning instructions are most prominent when available data (shots) are limited, emphasizing that optimal use of LLMs’ prior knowledge is vital for boosting performance in these contexts. Finally, our ensemble method yields consistent performance gains across all configurations.

\vspace*{-0.12in}
\paragraph{Training \& inference time comparison.} We evaluate and contrast the training and inference durations across various methods. All tests utilize the Adult dataset, employing a training set comprising 16 shots. By default, a single A100 GPU is employed, except for TabLLM, which utilizes four A100 GPUs to enable model parallelism. Detailed runtime statistics are displayed in Table~\ref{tab:cost}. It is evident that \textsf{FeatLLM} delivers notably efficient inference times, approaching those achieved by standard approaches. On the other hand, approaches leveraging LLMs, including In-context learning, TABLET, and TabLLM, incur much lengthier inference periods. Regarding the time required for training, \textsf{FeatLLM} demonstrates a similar scale to other LLM-driven methods—requiring just 30 API requests, which remains manageable from a cost perspective, and can further benefit from parallel execution.

\vspace*{-0.12in}
\paragraph{Quality of features generated by \textsf{FeatLLM}.}
To assess how informative the rule-based features produced by \textsf{FeatLLM} are in practical scenarios, we design a simple experiment. For each dataset, we calculate the mean AUROC for each generated feature relative to the target label, and then determine the proportion of features achieving an AUROC above 0.5. 
Results are reported in Table~\ref{tab:quality}. The findings indicate that the vast majority of these constructed rules hold meaningful relationships to the target variable and offer valuable signal for the downstream task (see the ``Before tuning" column). Further, during linear model training, features with weak associations to the response variable (those whose learned weights fall below a certain threshold) are automatically excluded (as detailed in the ``After tuning" column). For this analysis, the threshold is set at 0.1, reflecting the empirical distribution of the model's weight coefficients.

\vspace*{-0.12in}
\paragraph{Effect of adding spurious correlations.} To evaluate the capacity of different approaches to discard noisy, spurious correlations within a limited labeled data scenario, we performed a dedicated assessment. The Heart dataset served as the primary task dataset, into which we injected randomly selected columns from the Adult dataset that have no relevance to the Heart task. The impact of this perturbation was tracked by systematically increasing the number of these unrelated columns, and monitoring model performance. To maintain consistency, each method was given access to exactly 8 labeled samples, with no use of unlabeled data, thus omitting methods such as SCARF and STUNT from our set of baselines. 

The results are shown in Figure~\ref{fig:spurious}, which depicts performance shifts induced by introducing spurious correlations. Among all methods tested, \textsf{FeatLLM} exhibited the smallest reduction in performance as irrelevant features were added. We attribute this resilience to the model’s architecture, which leverages prior knowledge to explicitly align features with the specific task during rule extraction, making it less likely to generate rules involving noise features. This is corroborated by the observation that only 1–2\% of generated rules referenced these extraneous columns. Nevertheless, we also notice that \textsf{FeatLLM} is not entirely immune to erroneous reasoning: if columns sourced from data closely related to the main task are added, the model may erroneously internalize these spurious associations, as detailed in Appendix~\ref{sec:spurious-appendix}.

\vspace*{-0.12in}
\paragraph{Hyper-parameter analysis}
\textsf{FeatLLM} operates with two principal hyper-parameters: the ensemble size and the quantity of rules per class extracted in each LLM interaction. We performed a systematic evaluation of how these settings affect the model's performance. Our observations reveal that increasing the number of ensembles leads to better outcomes, but after reaching approximately 20 ensembles, these improvements plateau (refer to Fig.~\ref{fig:appendix-hyperparameter1} in Appendix). Regarding rule count, the performance differences are not statistically meaningful; nonetheless, we identify 10 as an optimal balance between prompt length and effectiveness (see Fig.~\ref{fig:hyperparameter2}). We conjecture this is because raising the number of rules increases the input’s dimensionality and the volume of information, but also tends to promote overfitting, particularly under low-shot conditions.

\section{Conclusion}
In this work, we introduce \textsf{FeatLLM}, which establishes new standards for LLM-based few-shot learning on tabular data. Rather than using LLMs merely for direct instance-level prediction, \textsf{FeatLLM} shifts the focus to creating decision criteria akin to those devised by feature engineers. By employing a rule generation strategy coupled with bagged ensemble modeling, \textsf{FeatLLM} minimizes inference overhead, functions purely with API-level access without any training, and addresses challenges posed by large feature spaces. The approach consistently delivers strong predictive performance and proves to be an effective, data-efficient alternative for leveraging LLMs on practical tabular datasets.

At present, \textsf{FeatLLM} targets scenarios typified by limited labeled data. Looking forward, our objective is to generalize the framework to handle richer datasets with more abundant samples, to consider generating diverse types of features beyond just rules, and to integrate interpretability mechanisms, broadening its utility across a spectrum of tasks.

\section{Broader Impact} 
The methodology underlying \textsf{FeatLLM} seeks to harness the pre-existing knowledge embedded in LLMs to elevate the performance ceiling for tabular learning, especially compared with conventional supervised approaches. By augmenting learning with this auxiliary knowledge, our work substantially enhances data efficiency and alleviates problems of overfitting stemming from small sample sizes. \textsf{FeatLLM} stands to benefit professionals in fields such as finance and healthcare—domains where manual labeling is costly or impractical and where LLMs pretrained on relevant textual data can deliver considerable prior knowledge advantages. A pressing direction for further study, critical to broader adoption, is to systematically examine how societal biases and misinformation present in LLM pretraining data might influence model predictions, as such prior knowledge could overrule patterns present in the available labeled data.

\end{document}